\section{Revision experiments}
\label{app:revision}

All runs in this appendix were added after review and pre-registered: predictions and
decision rules were written down and committed before any result was opened
(App.~\ref{app:prereg}), and selection uses synthetic validation only, as in the body.

\paragraph{Frozen-backbone block.} The frozen backbone runs in inference mode, so without
augmentation its pooled features (and the ImageNet logits computed from them) can be cached
once, and heads retrained on them in minutes. Applying the published seed-42 head to the
cache reproduces its published errors to within $0.0005$. Table~\ref{tab:frozen} gives three
seeds (42, 43, 44) for each configuration, all with the published head, loss and selection
rule. The 4{,}000-image cells use one fixed draw from the 28{,}000, not the original
4{,}000-image set. The tuned head is the best of learning rate $\{3{\times}10^{-4}, 10^{-3},
3{\times}10^{-3}\}\times\{20, 60\}$ epochs on seed 42, chosen on synthetic validation. The
synthetic-BN head re-estimates every BatchNorm layer on the un-augmented synthetic training
images, with the weights untouched, before caching.


\paragraph{Frozen backbone with augmentation.} Augmented inputs cannot use the feature cache,
so seeds 43 and 44 of the frozen augmented arm were trained end to end with the released code
on an Apple M5 GPU, while seed 42 is the published Colab run. Table~\ref{tab:frozenaug} pairs
each seed with the head-only retrain of the same seed in the published configuration
(Table~\ref{tab:frozen}), and their conformal coverage is in App.~\ref{app:conformal}.


\paragraph{Augmentation ablation.} Arm A is the full stack without blur and noise. Arm B keeps
only brightness and contrast jitter, gamma and random erasing (no glare blob, blur or noise).
Both train layer 4 on seeds 42, 43 and 44 with the published configuration, and we read them
as per-seed differences. They were trained with the released code on an Apple M5 GPU, whereas
the published arms were trained on Colab. A control rerun of the full stack (seed 42) on the
M5 differs from its published counterpart by $-0.053$, $+0.006$ and $-0.142$\,rad in $\Erot$
(synthetic, lightbox, sunlamp). That is within seed variation, but on sunlamp it is as large
as the ablation effects themselves.


\paragraph{Test-time BatchNorm adaptation (AdaBN).} For each layer-4 checkpoint we
re-estimate only layer 4's BatchNorm statistics, the ones fitted on synthetic data. We use
the unlabelled images of one half of each real domain (a fixed random split by filename),
leave every weight untouched, and evaluate on the other half (3{,}370 lightbox and 1{,}396
sunlamp images), with each arm's published synthetic-calibrated $\hat{q}$ for coverage. This is
test-time adaptation: it uses target images, though not their labels. The frozen arms are
excluded, as all their statistics are ImageNet's.


\paragraph{Image statistics.} Figure~\ref{fig:crops} compares pooled intensity densities,
which cannot tell a soft bright blob from a hard edge. Table~\ref{tab:imagestats} adds spatial
statistics, pooled over every image of each split after the same resize and grayscale
conversion (augmentation seed 7, as in Figure~\ref{fig:crops}):
\begin{itemize}[leftmargin=1.2em,itemsep=0pt,topsep=2pt]
  \item the Sobel gradient magnitude, scaled so a full black-to-white step reads 1
  \item edge density, the share of pixels whose gradient exceeds 0.1
  \item the clipped fraction, the share of pixels at or above $254/255$
  \item the total variation (TV) between pooled densities (48 intensity bins, 64 gradient
  bins)
\end{itemize}
The intensity TV of the full stack reproduces Figure~\ref{fig:crops} ($0.63$, $0.80$).
